\section{Continuous Assessment 2 --- Set A}
\label{sec:ca2_set_a}

\ExamHeader{Continuous Assessment 2}{A}{45 Minutes}{25 Marks}{Unit 4 (Number Systems) \& Unit 5 (Version Control)}

\noindent
\textbf{Instructions:} All 25 questions are compulsory. Each question carries 1 mark. Select the single best answer.

\vspace{3mm}

\mcqitem{1}{What is the base (radix) of the Binary number system?}{2}{8}{10}{16}

\mcqitem{2}{What is the base (radix) of the Octal number system?}{2}{8}{10}{16}

\mcqitem{3}{What is the base (radix) of the Hexadecimal number system?}{8}{10}{12}{16}

\mcqitem{4}{Which of the following is an INVALID octal number?}{107}{543}{782}{615}

\mcqitem{5}{In the hexadecimal system, what decimal value is represented by the letter 'A'?}{10}{11}{12}{15}

\mcqitem{6}{In the hexadecimal system, what decimal value is represented by the letter 'F'?}{10}{12}{14}{15}

\mcqitem{7}{What is the decimal equivalent of the binary number $(1011)_2$?}{9}{11}{13}{15}

\mcqitem{8}{How many binary bits correspond directly to a single Octal digit?}{2 bits}{3 bits}{4 bits}{8 bits}

\mcqitem{9}{How many binary bits correspond directly to a single Hexadecimal digit?}{2 bits}{3 bits}{4 bits}{8 bits}

\mcqitem{10}{What is the octal equivalent of the binary number $(110101)_2$?}{$(55)_8$}{$(65)_8$}{$(35)_8$}{$(71)_8$}

\mcqitem{11}{What is the hexadecimal equivalent of the binary number $(10101111)_2$?}{$(\text{AF})_{16}$}{$(\text{FA})_{16}$}{$(\text{BF})_{16}$}{$(\text{AE})_{16}$}

\mcqitem{12}{What is the decimal value of the binary fraction $(0.1)_2$?}{0.1}{0.5}{0.25}{0.125}

\mcqitem{13}{In a positional number, what is the leftmost non-zero digit called?}{Least Significant Digit (LSD)}{Most Significant Digit (MSD)}{Radix digit}{Fractional bit}

\mcqitem{14}{Who designed and created the Git version control system in 2005?}{Dennis Ritchie}{Linus Torvalds}{James Gosling}{Guido van Rossum}

\mcqitem{15}{What type of version control architecture is used by Git?}{Centralized only}{Distributed version control system}{Peer-to-peer torrent system}{Manual logging}

\mcqitem{16}{What is GitHub?}{A command-line text editor}{A cloud-based hosting and collaboration service for Git repositories}{A programming language compiler}{A CPU component}

\mcqitem{17}{Which hidden directory is created by Git inside a project folder to store tracking metadata?}{\texttt{.github}}{\texttt{.git}}{\texttt{.version}}{\texttt{.repo}}

\mcqitem{18}{Which command initializes a new local Git repository in the current directory?}{\texttt{git start}}{\texttt{git init}}{\texttt{git create}}{\texttt{git new}}

\mcqitem{19}{Which Git command checks the status of modified, staged, and untracked files?}{\texttt{git status}}{\texttt{git check}}{\texttt{git log}}{\texttt{git inspect}}

\mcqitem{20}{In Git's three-stage architecture, what is the intermediate area where changes are prepared before being committed?}{Working Directory}{Staging Area (Index)}{Remote Server}{Trash Bin}

\mcqitem{21}{Which command is used to move changes from the working directory to the Git staging area?}{\texttt{git commit}}{\texttt{git add <filename>}}{\texttt{git save}}{\texttt{git stage}}

\mcqitem{22}{Which command permanently records the staged snapshot into the local repository history?}{\texttt{git save}}{\texttt{git commit -m "message"}}{\texttt{git push}}{\texttt{git freeze}}

\mcqitem{23}{What is the purpose of the \texttt{.gitignore} file?}{To list files that Git should intentionally ignore and not track}{To delete old commits}{To store passwords}{To speed up downloads}

\mcqitem{24}{In Git, what is a branch?}{A separate hard drive}{A lightweight movable pointer to a specific commit}{A deleted repository}{An untracked image file}

\mcqitem{25}{Which command uploads local branch commits to a remote repository like GitHub?}{\texttt{git pull}}{\texttt{git push}}{\texttt{git fetch}}{\texttt{git commit}}

\newpage
\subsection*{Answer Key \& Explanations --- CA-2 (Set A)}

\begin{answerkeytable}{CA-2 Set A --- Answer Key}
\begin{tabular}{*{10}{c}}
\toprule
\textbf{Q} & \textbf{Ans} & \textbf{Q} & \textbf{Ans} & \textbf{Q} & \textbf{Ans} & \textbf{Q} & \textbf{Ans} & \textbf{Q} & \textbf{Ans} \\
\midrule
1 & \textbf{A} & 6 & \textbf{D} & 11 & \textbf{A} & 16 & \textbf{B} & 21 & \textbf{B} \\
2 & \textbf{B} & 7 & \textbf{B} & 12 & \textbf{B} & 17 & \textbf{B} & 22 & \textbf{B} \\
3 & \textbf{D} & 8 & \textbf{B} & 13 & \textbf{B} & 18 & \textbf{B} & 23 & \textbf{A} \\
4 & \textbf{C} & 9 & \textbf{C} & 14 & \textbf{B} & 19 & \textbf{A} & 24 & \textbf{B} \\
5 & \textbf{A} & 10 & \textbf{B} & 15 & \textbf{B} & 20 & \textbf{B} & 25 & \textbf{B} \\
\bottomrule
\end{tabular}
\end{answerkeytable}

\begin{expbox}[Technical Solutions --- CA-2 Set A]
\begin{itemize}[leftmargin=5mm]
    \item \textbf{Q.4 (C):} The digit '$8$' is invalid in Octal (base 8 allows digits $0$ to $7$).
    \item \textbf{Q.7 (B):} $1011_2 = 1(8) + 0(4) + 1(2) + 1(1) = 11_{10}$.
    \item \textbf{Q.10 (B):} $110101_2$ grouped in 3-bit sets: $\underbrace{110}_{6}\;\underbrace{101}_{5} = (65)_8$.
    \item \textbf{Q.11 (A):} $10101111_2$ grouped in 4-bit sets: $\underbrace{1010}_{\text{A}}\;\underbrace{1111}_{\text{F}} = (\text{AF})_{16}$.
    \item \textbf{Q.14 (B):} Linus Torvalds created Git in 2005 to manage Linux kernel source code.
    \item \textbf{Q.21 (B):} \texttt{git add} moves files into the staging area (index).
    \item \textbf{Q.25 (B):} \texttt{git push} transmits commits from the local branch to the remote repository.
\end{itemize}
\end{expbox}

\newpage
